\documentclass{jsarticle}
\usepackage{amsmath,amssymb, stmaryrd}
\usepackage[dvipdfmx]{graphicx}
\begin{document}
\title{数学特別講義4 Leanコード課題}
\author{linx03642-cyber}
\date{\today}
\maketitle

本課題では, 次の問題に対してLeanによる形式証明を与えた.

\begin{quote}
定数記号$B, C, I$と2項関数記号$f$から生成される項全体の集合を$\mathcal{T}$とする(変数は含まないものとする).
つまり$\mathcal{T}$は$B, C, I$を含み, $M, N \in \mathcal{T}$ならば$f(M, N) \in \mathcal{T}$を満たす最小の集合である.

同様にして, 定数記号0と2項関数記号$g$から生成される項全体の集合を$\mathcal{U}$とする(同じく変数は含まないものとする).
以降$g(a, b)$のことを$(a \to b)$と書く.
また$(a \to (b \to c))$を$(a \to b \to c)$と略記し, $(a \to (b \to (c \to d)))$を$(a \to b \to c \to d)$と略記する.

各項$M_{0} \in \mathcal{T}$に対して集合$\llbracket M_{0} \rrbracket \subseteq \mathcal{U}$を次のように帰納的に定める.
\begin{align*}
    \llbracket B \rrbracket &= \{((b \to c) \to (a \to b) \to (a \to c)) \mid a, b, c \in \mathcal{U}\} \\
    \llbracket C \rrbracket &= \{((a \to b \to c) \to (b \to a \to c)) \mid a, b, c \in \mathcal{U}\} \\
    \llbracket I \rrbracket &= \{(a \to a) \mid a \in \mathcal{U}\} \\
    \llbracket f(M, N) \rrbracket &= \{b \mid あるa \in \mathcal{U} について(a \to b) \in \llbracket M \rrbracket, a \in \llbracket N \rrbracket\}.
\end{align*}

以下の各集合$X_{i} (i = 1, 2, 3, 4)$について, $X_{i} = \llbracket M_{i} \rrbracket$となる項$M_{i} \in \mathcal{T}$は存在するかどうか, 理由をつけて答えよ.
\begin{align*}
    X_{1} &= \{(0 \to 0)\} \\
    X_{2} &= \{(a \to b \to b \to a) \mid a, b \in \mathcal{U}\} \\
    X_{3} &= \{(((a \to b) \to b) \to a) \mid a, b \in \mathcal{U}\} \\
    X_{4} &= \{((a \to b \to c \to d) \to (c \to b \to a \to d)) \mid a, b, c, d \in \mathcal{U}\}.
\end{align*}

(出典: 平成29年度 京都大学大学院理学研究科数学・数理解析専攻 数学系・数理解析系入学試験問題 専門科目 問10)
\end{quote}

本問は大学院入試の過去問であり, 公式の解答は公開されていない.
そのため, Leanを用いて形式化を行うことで自身の証明の正当性を確認することは, 自らの学習において有意義であるほか, 本授業の主題である「数学系エンドユーザーのためのLean入門」とも重なるため, 本課題として取り組むこととした.
特に本問は, 計算機科学や数理論理学の領域に属し, 自身の興味とも合致する.
問題の構成としても, それぞれの$X_{i}$に対して別々のアプローチが求められ, 構成的な存在証明も非存在証明も含まれるため, 授業内で学習したLeanの記法や証明手法を広く活用できると考え, 本問の形式的証明を与えた.

「存在しない」が答えである$X_{1}, X_{2}, X_{3}$について, 紙の上で自力で問題を解いた際には全く異なるアプローチをとっているように感じていた.
$X_{1}$では, 直感的に有限集合が生成されるとは考えづらく, 具体的にある元が存在するなら, それへの代入によって得られる項も含まれるという発想に至る.
また, $X_{3}$では, 意味論と対応させて, 生成される項がトートロジーであることを示せばよいため, 自然に思いつきやすい.
一方, $X_{2}$では, 各項の出現の正負の数を数えることで, それらの出現が同数であることが保存されているという線形論理的な発想が求められ, ひらめくことが難しかった.

しかし, これらの発想は形式化を行う中で, ほとんど同じような構造を持った証明であることに気づき, メタ的な視点からこれらの問題を捉えることができた.
形式化を行う中で最も面白さを感じたのは, このように証明の本質的な構造を捉えることができた点である.
いずれの問題も(広い意味での)不変量を構成し, それが確かに不変量であることを構造帰納法で示し, 最後に反例を構成して不変性に反することを示すことによって解決した.
一方, 普段の数学では自明とする点に関しても, どの型でどのように自明であるかを明示しなければならない点は, 形式化を行う上での難しさであるように感じた.

その他の工夫として, 本課題の作成にあたっては, 強力な自動化タクティクの使用を控え, 基礎的なタクティクを中心にして証明を構成した.
これは, 初めてLeanで証明を書くにあたり, 論理のステップやCurry-Howard同型対応を意識しながら証明を構成することがLeanの理解において有意義であると考えたためである.
\end{document}